% The ACM Web Conference 2027, Research Track, Semantics and Knowledge.
% https://www2027.thewebconf.org/research-track-papers/
% Submission format: \documentclass[sigconf, review, anonymous]{acmart};
% at most 8 pages of content plus up to 4 pages of references and appendix.
\documentclass[sigconf, review, anonymous]{acmart}
% \documentclass[sigconf]{acmart} % camera-ready

% acmart already loads amsmath, booktabs, graphicx, and newtxmath;
% do not load amssymb, which clashes with newtxmath.
\usepackage{multirow}
\usepackage{array}
\usepackage{pifont}
\newcolumntype{C}[1]{>{\centering\arraybackslash}p{#1}} % fixed-width centred column for results tables
\setlength{\emergencystretch}{2em} % fewer overfull lines with long system names

\usepackage{todonotes}
% \usepackage[disable]{todonotes}
\newcommand{\oshani}[1]{\todo[inline, color=blue!40, author=\textbf{Oshani}]{#1}}
\newcommand{\maruf}[1]{\todo[inline, color=red!40, author=\textbf{Maruf}]{#1}}
\newcommand{\abid}[1]{\todo[inline, color=green!40, author=\textbf{Abid}]{#1}}

\settopmatter{printacmref=false}
\setcopyright{none}
\renewcommand\footnotetextcopyrightpermission[1]{}

% ---------------------------------------------------------------------------
% Notation. \sys is the system name everywhere in the text; change it in one
% place. \tbd marks a result cell that is filled after the UNH runs finish.
% ---------------------------------------------------------------------------
\newcommand{\sys}{UNH}
\newcommand{\sysfull}{Unified Narrative Hypergraph}
\newcommand{\tbd}{--}
\providecommand{\argmax}{\operatorname*{arg\,max}}
\newcommand{\rel}[1]{\textsc{#1}}            % relation label, e.g. \rel{causes}
\newcommand{\Ecf}{E_{CF}}
\newcommand{\Efv}{E_{FV}}
\newcommand{\Eff}{E_{FF}}
\newcommand{\simcos}{\operatorname{cos}}
\newcommand{\best}{\operatorname{best}}
\newcommand{\cover}{\operatorname{cover}}
% feature table marks
\newcommand{\yes}{\ding{51}}
\newcommand{\no}{\ding{55}}
\newcommand{\half}{$\circ$}
\newcommand{\na}{--}
% figure that compiles before its file exists
\newcommand{\figorbox}[3]{\IfFileExists{#1}{\includegraphics[width=#2]{#1}}{\fbox{\parbox[c][#3][c]{0.95#2}{\centering\small Placeholder: \texttt{\detokenize{#1}}}}}}

\begin{document}

\title{Facts as Hyperedges and Discourse Nodes: A Unified Narrative Hypergraph for Graph Retrieval-Augmented Generation}

\author{Abid Talukder}
\affiliation{%
  \institution{Rensselaer Polytechnic Institute}
  \city{Troy}
  \state{New York}
  \country{USA}}
\email{taluka@rpi.edu}

\author{Maruf Ahmed Mridul}
\affiliation{%
  \institution{Rensselaer Polytechnic Institute}
  \city{Troy}
  \state{New York}
  \country{USA}}
\email{mridum@rpi.edu}

\author{Oshani Seneviratne}
\affiliation{%
  \institution{Rensselaer Polytechnic Institute}
  \city{Troy}
  \state{New York}
  \country{USA}}
\email{senevo@rpi.edu}

\renewcommand{\shortauthors}{Talukder et al.}

\begin{abstract}
Graph-based retrieval-augmented generation (RAG) has a large language model (LLM) turn a corpus into a knowledge graph and retrieves along the graph's links to ground the LLM's answers. In existing indices, two pieces of knowledge are connected in one of two ways. Either they mention the same entity, whether the index stores triples, $n$-ary facts, or propositions, or the LLM fuses them into a summary. Many questions instead follow relations between \emph{statements}. Consider whether a life-insurance policy pays out after the insured was disabled in a drunk-driving crash. The answer depends on a chain of clauses from different sections: an exception that removes the premium waiver, a grace period, a lapse rule, and the death benefit. These clauses share only the words ``premium'' and ``policy'', which appear in 24\% and 35\% of the policy's sentences, so an entity hop cannot tell which clause comes next. A community summary can relate only the clauses its cluster happens to contain, and it paraphrases the very conditions that decide the answer. We present the \sysfull{} (\sys{}), an index in which every atomic fact extracted from the corpus is both a hyperedge over the entities that take part in it and a node in a graph of typed relations between facts, such as \rel{exception\_to}, \rel{leads\_to}, and \rel{precedes}. An LLM decides these relations, each is stored as a label rather than as text, and every fact stays linked to the passage that states it. To find relations across documents without comparing every pair, \sys{} matches the follow-up questions each fact raises against the questions other facts answer. It clusters facts by this affinity and judges each cluster in one LLM call, so the number of calls grows linearly with the corpus. At query time, one LLM call predicts which relation types the answer should follow. A beam search prefers those relations to generic entity hops, and the reader receives short chains of source facts joined by their relation labels. We compare \sys{} with six graph-based RAG systems on HotpotQA, MuSiQue, and QuALITY under a protocol that fixes the LLM, embedding model, reader, and judge. We isolate the relation graph by evaluating hypergraph-only and relation-only views of the same index, and we test summaries written over the same clusters.
\abid{Add the headline result sentence once the UNH runs are scored.}
\end{abstract}

\begin{CCSXML}
<ccs2012>
   <concept>
       <concept_id>10002951.10003317.10003338</concept_id>
       <concept_desc>Information systems~Retrieval models and ranking</concept_desc>
       <concept_significance>500</concept_significance>
   </concept>
   <concept>
       <concept_id>10010147.10010178.10010187</concept_id>
       <concept_desc>Computing methodologies~Knowledge representation and reasoning</concept_desc>
       <concept_significance>500</concept_significance>
   </concept>
   <concept>
       <concept_id>10010147.10010178.10010179.10003352</concept_id>
       <concept_desc>Computing methodologies~Information extraction</concept_desc>
       <concept_significance>300</concept_significance>
   </concept>
</ccs2012>
\end{CCSXML}

\ccsdesc[500]{Information systems~Retrieval models and ranking}
\ccsdesc[500]{Computing methodologies~Knowledge representation and reasoning}
\ccsdesc[300]{Computing methodologies~Information extraction}

\keywords{Retrieval-augmented generation, knowledge graph construction, knowledge hypergraphs, discourse relations, multi-hop question answering}

\maketitle

% =============================================================================
\section{Introduction}
% =============================================================================

Retrieval-augmented generation (RAG) grounds a large language model (LLM) in an external corpus~\cite{lewis2020retrieval}. Graph-based RAG goes one step further. An LLM first turns the corpus into a knowledge graph, and retrieval then follows the graph's links to collect evidence that is related to other evidence, not only to the question~\cite{peng2025graph,han2024retrieval,edge2024local,gutierrez2025rag}. These graphs are a new kind of knowledge graph on the Web: built automatically from text, specific to one corpus, and read by a machine. What such a graph can connect determines what the machine reading it can be told.

\begin{figure*}[t]
  \centering
  \figorbox{figures/motivation.pdf}{\textwidth}{5.6cm}
  \caption{\emph{Is the death benefit paid?} The answer follows five clauses from four sections of a term-life policy: the premium waiver ($f_1$), its drunk-driving exception ($f_2$), the grace period ($f_3$), the lapse rule ($f_4$), and the death benefit ($f_5$). (a)~Dense retrieval returns the clauses that share the question's words and misses the grace period and the lapse rule. (b)~A triple graph and (c)~a hypergraph connect $f_2$ to $f_3$, and $f_4$ to $f_5$, only through \emph{premium} and \emph{policy}, entities that dozens of other clauses also mention. (d)~A community summary relates only the clauses in its own cluster. The exception sits with the other disability clauses and the lapse rule with the other payment clauses, so the link the question needs is written in neither summary. (e)~\sys{} links the clauses by typed relations, so the answer is a path in the index: $f_2$ is an \rel{exception\_to} $f_1$ and \rel{leads\_to} $f_3$, which \rel{precedes} $f_4$, which \rel{prevents} $f_5$.}
  \label{fig:motivation}
\end{figure*}

\paragraph{A motivating example.}
Insurance policies are published on the Web, and insured people, or agents acting for them, ask coverage questions about them. Consider this question about a term-life policy (Figure~\ref{fig:motivation}): ``\emph{The insured was totally disabled in a crash while driving over the blood-alcohol limit. She relied on the policy's waiver of premiums and died a year later. Is the death benefit paid?}'' The answer is no, and it follows five clauses from four sections of the policy.\footnote{Clauses abridged from a public sample of a 20-year term-life policy.}
\begin{itemize}
  \item The waiver section states that premiums are waived while the insured is totally disabled ($f_1$).
  \item Its list of exceptions states that premiums are not waived if the disability results from an accident in which the insured's blood alcohol exceeded the legal limit ($f_2$).
  \item The payment section, pages earlier, allows a 30-day grace period for unpaid premiums ($f_3$). It also states that the policy lapses if they remain unpaid at the end of that period ($f_4$).
  \item The policy schedule states that the sum insured is paid if the insured dies while the policy is in effect ($f_5$).
\end{itemize}
The exception removed the waiver, so premiums remained due. Left unpaid, they let the policy lapse, and a lapsed policy pays nothing. No clause states this chain. Its middle links, the grace period and the lapse rule, share no word with the question other than \emph{premium}.
\abid{Add the source of the sample policy (TD 20-Year Term Life sample policy package) in the camera-ready version; keep it generic for review.}

\paragraph{Why entity links are not enough.}
The clauses are related by exception, consequence, and condition, not by what they mention. The exception speaks of disability and drunk driving, the grace period of due dates, and the death benefit of the beneficiary. The only entities they share are \emph{premium} and \emph{policy}, which occur in 60 and 88 of the policy's 251 sentences. Each family of indices fails on this chain in its own way (Figure~\ref{fig:motivation}a--c). Dense retrieval returns the clauses most similar to the question (the waiver, the exception, and the death benefit) and misses the grace period and the lapse rule that decide the answer. A triple graph~\cite{edge2024local,gutierrez2025rag} and a hypergraph of $n$-ary facts~\cite{luo2025hypergraphrag} connect the exception to the grace period only through the node \emph{premium}. There the grace period is one neighbour among sixty, and nothing on the edges says that a denied waiver leads to unpaid premiums or that unpaid premiums lead to a lapse. The hop exists, but it carries no information about which neighbour comes next. Where the dependent clauses share no entity at all, there is no hop.

\paragraph{Why not community summaries?}
A natural remedy is to let the LLM write the connections down. Microsoft GraphRAG~\cite{edge2024local}, RAPTOR~\cite{sarthi2024raptor}, and HyGRAG~\cite{zhong2026hygrag} cluster the corpus or its graph and write a summary of each cluster. A summary can fuse what several passages say and state what none of them states alone. One could even prompt the summaries to spell out every causal, temporal, and conditional relation among their members. Three problems remain, and a better prompt solves none of them.
\begin{itemize}
  \item \emph{A summary can only relate what its cluster contains.} Clusters are formed from shared entities or from embedding similarity. The drunk-driving exception therefore falls with the other disability clauses, and the lapse rule with the other payment clauses. The relation the question needs lies \emph{between} the two clusters and is written in neither summary.
  \item \emph{A summary is fixed before any question is asked.} It must choose which of the many relations among its members to state, at one granularity, in a paragraph the reader receives whole. It cannot be cut down to the one chain a question needs, or chained with another cluster's summary.
  \item \emph{A summary paraphrases.} In a contract, the scope of a clause is often the answer. The exception removes the premium waiver, not the coverage. A summary that reads ``disabilities caused by drunk driving are excluded'' gives the wrong answer for an insured who kept paying.
\end{itemize}
Summaries remain the right tool for questions about a corpus as a whole~\cite{edge2024local}. For questions that follow a chain of specific statements, the connections must be stored where retrieval can select them one at a time, between facts that the reader can check against their source. The fusion then happens at query time, by the reader, over exactly the facts the question needs.

\paragraph{The gap.}
Existing indices therefore have no way to represent the relation between two statements: it is absent, reduced to a hub entity, or dissolved into generated text. This has two consequences. First, \emph{dependent facts that share no informative entity cannot be found from each other}, although causes and effects, rules and their exceptions, and claims and their rebuttals are routinely stated with different terms and in different places. The same pattern links a cause reported by one news outlet to an effect reported by another, a finding to the result that overturned it, and, in long narratives, an event to the motive behind it~\cite{pang2022quality}. Second, \emph{a step from one fact to the next says nothing about how they relate}, so retrieval cannot prefer the consequence of an exception to the dozens of other facts about premiums. The Semantic Web has long treated statements as objects in their own right. $N$-ary relation patterns~\cite{noy2006nary}, Wikidata's qualified statements~\cite{vrandecic2014wikidata}, RDF-star~\cite{hartig2017rdfstar}, and nanopublications~\cite{groth2010nanopub} all give a statement an identity to which further assertions can attach. Graph RAG indices have adopted half of this idea by making $n$-ary facts the unit of retrieval~\cite{luo2025hypergraphrag,wang2025proprag}. They have not adopted the other half: facts are never related to one another.
\abid{Strengthen with one number: the share of gold reasoning chains in QuALITY and MuSiQue whose consecutive supporting facts share no entity, or share only an entity mentioned by more than 5\% of facts, e.g., measured on 100 annotated questions.}

\paragraph{Our approach.}
We introduce the \textbf{\sysfull{} (\sys{})}, an index in which the relation between two facts is itself an element of the graph. Every atomic fact extracted from the corpus plays two roles. It is a \emph{hyperedge} over the entities that take part in it. It is also a \emph{node} in a directed graph of relations between facts, where each relation has a short type label such as \rel{exception\_to}, \rel{leads\_to}, \rel{precedes}, or \rel{contradicts}. Every fact keeps a link to the passage that states it. In Figure~\ref{fig:motivation}e, the exception $f_2$ is an \rel{exception\_to} the waiver $f_1$ and \rel{leads\_to} the grace period $f_3$. The grace period \rel{precedes} the lapse rule $f_4$, which \rel{prevents} the death benefit $f_5$. The sixty other facts about premiums remain attached to the entity \emph{premium} but not to this chain. A retrieval path can therefore enter at the clause that matches the question, follow typed relations across sections, and reach clauses that share none of the question's words. This path exists in neither the hypergraph nor the relation graph alone. The reader receives this chain, with each clause traceable to its section, and draws the conclusion itself.

Two problems stand between this representation and a practical index. The first is \emph{which facts to relate}. Asking an LLM about every pair of facts is quadratic in corpus size. Embedding similarity alone misses exactly the pairs we want, because the exception and the grace period have little in common. \sys{} annotates each fact with the questions it answers and the follow-up questions it raises. It then scores a candidate pair by how well the questions one fact raises (\emph{what happens if premiums are not waived and go unpaid?}) match those the other answers (\emph{what happens when a premium is not paid by its due date?}). Facts are clustered by this affinity, and one LLM call decides the relations within each cluster, so the number of calls grows linearly with the corpus. Unlike a summary, the result of the call is a set of labelled edges. Edges found in different clusters can be chained by retrieval, and each edge can be selected on its own. The second problem is \emph{how retrieval should use the relations}. A single LLM call turns the question into a plan that names the relation types the answer path should follow. A beam search prefers edges whose labels match the plan. The reader then receives short, deduplicated chains of source facts joined by their relation labels.

We compare \sys{} with six graph-based RAG systems: HippoRAG~2~\cite{gutierrez2025rag}, LogicRAG~\cite{chen2026logicrag}, Mem\-Graph\-RAG~\cite{wu2026memgraphrag}, Hyper\-Graph\-RAG~\cite{luo2025hypergraphrag}, Microsoft Graph\-RAG~\cite{edge2024local}, and Hy\-GRAG~\cite{zhong2026hygrag}. The benchmarks are multi-hop questions over Wikipedia (HotpotQA~\cite{yang2018hotpotqa}, MuSiQue~\cite{trivedi2022musique}) and long narratives (QuALITY~\cite{pang2022quality}). Every system uses the same LLM, embedding model, reader prompt, and judge, so differences come from indexing and retrieval alone. Our contributions are:
\begin{itemize}
  \item \textbf{Representation.} \sys{}, a knowledge-graph model for RAG in which atomic facts are both $n$-ary hyperedges over entities and nodes of a graph of typed relations, with passage-level provenance. Hypergraph and relation-graph indices are edge-restricted views of it (Section~\ref{sec:index}).
  \item \textbf{Construction.} A pipeline that finds typed relations between facts across documents with a number of LLM calls linear in corpus size, using question-based affinity and community-level judgment that outputs edges rather than summaries (Section~\ref{sec:construction}).
  \item \textbf{Retrieval.} A procedure with one LLM call per question that seeds from facts and entities jointly, follows the relation types predicted from the question, and returns grounded fact chains (Section~\ref{sec:retrieval}).
  \item \textbf{Evaluation.} A controlled comparison of seven systems that reports accuracy together with construction tokens and reader context. It includes ablations over views of the same index, which isolate the relation graph, and a variant that writes relation-focused summaries over the same clusters instead of edges (Section~\ref{sec:experiments}).
\end{itemize}
\abid{Add one sentence with the headline findings after the UNH runs.}
\abid{Verify panel (d) of Figure~\ref{fig:motivation} on a real index: build GraphRAG (or HyGRAG) on the sample policy and check that the drunk-driving exception and the lapse rule land in different communities. Quote the actual summary text in the figure.}

% =============================================================================
\section{Related Work (Claude Version)}
\label{sec:related_claude}
% =============================================================================
% Second version of the related-work section. The original section above is
% kept unchanged; keep whichever version (or merge) before submission.

We group prior work by what its index can connect, because that determines which questions retrieval can follow. For each group we state what it does and what it leaves out. Table~\ref{tab:positioning} summarizes the comparison as capabilities.

\paragraph{Entity-centric graph RAG.}
The dominant design has an LLM extract entity--relation triples from each passage and merge them into one knowledge graph~\cite{bian2025llm,zhang2024extract,mo2025kggen,lairgi2024itext2kg}. Microsoft GraphRAG~\cite{edge2024local} partitions the graph into communities and pre-generates a report for each. LightRAG~\cite{guo2024lightrag} retrieves entities and relations through keywords. HippoRAG~\cite{gutierrez2024hipporag} and HippoRAG~2~\cite{gutierrez2025rag} run Personalized PageRank from the question's entities, and HippoRAG~2 adds passage nodes and LLM filtering of retrieved triples. MemGraphRAG~\cite{wu2026memgraphrag} shares a global memory between extraction agents to keep the schema consistent and to resolve conflicting triples. Iterative reasoners~\cite{sun2023think,ma2024think,luo2023reasoning} and learned retrievers~\cite{mavromatis2024gnn} walk the same triples. Systematic comparisons find that the graph pays off mainly when a question requires connecting evidence~\cite{zhou2025indepth,xiang2025when}. These methods leave out two things. A fact with more than two participants is split into triples that no longer record which participants belong to the same event. Facts are connected only through entity nodes, so a hub entity offers all of its facts as equal neighbours.

\paragraph{$N$-ary facts and propositions.}
A second line keeps a fact's participants together. HyperGraphRAG~\cite{luo2025hypergraphrag} stores each $n$-ary fact as a hyperedge over its entities. HyperRAG~\cite{lien2026hyperrag} learns to reason over such facts, HOLMES~\cite{panda2024holmes} builds a hyper-relational graph per query, and Cog-RAG~\cite{hu2026cograg} pairs an entity hypergraph with a hypergraph of chunk themes. Other methods change the retrieval unit to the atomic proposition~\cite{chen2024dense,min2023factscore}. GraphReader~\cite{li2024graphreader} attaches atomic facts to key elements of a long document. PropRAG~\cite{wang2025proprag} treats propositions as hyperedges and beam-searches over proposition paths. SiReRAG~\cite{zhang2025sirerag} groups propositions that share entities and summarizes each group. AtomicRAG~\cite{hou2026atomicrag} and ToPG~\cite{delmas2026topg} link atomic facts to entities and passages. These methods recover \emph{who} takes part in a fact. The connection between two facts, however, is still co-membership of an entity, a theme, or a chunk. It is untyped, and it does not exist when two facts share none.

\paragraph{Events, causality, and discourse.}
The methods closest to ours relate statements directly. ASER~\cite{zhang2020aser}, presented at this conference, built an eventuality graph whose edges are discourse relations of the Penn Discourse Treebank~\cite{prasad2017penn}, mined from explicit connectives. It is a commonsense resource, not an index of a given corpus. EventRAG~\cite{yang2025eventrag} extracts events with typed relations such as \rel{cause} and \rel{precede}. Its relations come from the extraction call of one chunk, under-connected events are expanded with links from shared participants or from the LLM's own knowledge, and an agent retrieves in several LLM rounds. E$^2$RAG~\cite{zhang2026e2rag} and DyG-RAG~\cite{sun2025dygrag} link events by co-mention, shared entities, or temporal proximity, without relation types. CausalRAG~\cite{wang2025causalrag} builds causal edges between concepts and asks an LLM to find causal paths at query time. Disco-RAG~\cite{liu2026discorag} parses rhetorical relations within and between the passages retrieved for a question and uses them to plan the answer. These methods leave out typed relations that are (i)~found across documents, (ii)~between grounded statements rather than concepts or event names, (iii)~built once over the whole corpus, and (iv)~used to decide \emph{what} to retrieve. A query-time structure over already retrieved passages cannot recover evidence that similarity search missed.

\paragraph{Generated summaries and query-time structure.}
Some indices fuse knowledge across passages by writing new text. RAPTOR~\cite{sarthi2024raptor} recursively summarizes clusters of chunks, NodeRAG~\cite{xu2025noderag} adds summary nodes to a heterogeneous graph, and HyGRAG~\cite{zhong2026hygrag} clusters a hybrid chunk--entity graph and writes summaries that fuse context and relations. Agent memories~\cite{xu2025amem,qian2024memorag} link or rewrite stored notes with an LLM. Other methods build structure per question. IRCoT~\cite{trivedi2023interleaving} interleaves retrieval with reasoning, and Plan$^{*}$RAG~\cite{verma2025planrag} and LogicRAG~\cite{chen2026logicrag} decompose the question into a graph of sub-questions; LogicRAG also keeps a rolling LLM summary of the evidence. Summaries can state what no single passage states, which suits questions about a corpus as a whole~\cite{edge2024local}. For questions that follow a chain of statements, they have three limits regardless of prompt (Section~1). They relate only what their cluster contains, they fix at index time which relations to state, and they paraphrase conditions the reader cannot trace to a source. Section~\ref{sec:ablation} tests the strongest version of this alternative: relation-focused summaries written over \sys{}'s own clusters. Query-time structure is discarded after each question, and its LLM cost recurs on every question (Section~\ref{sec:efficiency}).

\paragraph{Statements as knowledge objects.}
Knowledge representation on the Web has long given statements an identity. The W3C pattern for $n$-ary relations~\cite{noy2006nary} introduces a node for each relation instance, Wikidata attaches qualifiers and references to statements~\cite{vrandecic2014wikidata}, RDF-star lets a triple be the subject of another triple~\cite{hartig2017rdfstar}, and nanopublications bundle an assertion with its provenance~\cite{groth2010nanopub}. Discourse theory supplies the relations between statements: cause, contrast, elaboration, and temporal succession~\cite{mann1988rhetorical,prasad2017penn}. \sys{} applies both ideas to indices that an LLM builds from text. A fact is a statement node with its participants (the $n$-ary pattern) and its source passage (its provenance), and typed relations hold between facts. We do not fix a relation inventory; labels are open-vocabulary, produced under constraints (Section~\ref{sec:construction}), and may join facts from different documents.

\paragraph{The gaps \sys{} closes.}
Table~\ref{tab:positioning} distills four gaps. \textbf{(G1)~Connectivity:} entity-centric and $n$-ary indices connect facts only through shared entities, and the methods with relations between statements find them within one chunk or among already retrieved passages. \sys{} finds typed relations across the whole corpus, including between facts that share no entity. \textbf{(G2)~Guidance:} no corpus-level index chooses its next hop by the kind of relation a question asks for. \sys{} predicts relation types from the question and follows edges with matching labels. \textbf{(G3)~Grounding:} methods that connect distant evidence do so through generated summaries or the LLM's own knowledge, which fix the connections before the question is known and paraphrase the facts they connect. \sys{} stores each connection as a labelled edge that retrieval selects per question, and gives the reader only extracted facts, each traceable to its passage. \textbf{(G4)~Cost:} methods that build structure per question pay for it on every question. \sys{} builds its structure once with a number of LLM calls linear in the corpus and needs one LLM call per question.

\begin{table*}[t]
\centering
\caption{Capabilities of graph-based RAG methods. \textbf{$N$-ary facts}: one element of the index holds all participants of a fact. \textbf{Typed fact relations}: the index stores relations between facts (or events or passages), labelled by the kind of dependence. \textbf{No-shared-entity links}: the index can connect two facts from different passages that share no entity. \textbf{Relation-guided hops}: retrieval chooses its next step by the type of relation the question asks for. \textbf{Grounded context}: the reader sees only text from the corpus, each piece traceable to its source; relation labels count as types, not text. \textbf{Single-pass query}: at most one LLM call before the reader, over an index built once. \yes{}: yes; \half{}: partly (see note); \no{}: no; \na{}: not applicable, because the method builds no corpus index.}
\label{tab:positioning}
\small
\setlength{\tabcolsep}{4pt}
\begin{tabular}{l*{6}{C{1.9cm}}}
\toprule
 & \multicolumn{3}{c}{\textbf{What the index connects}} & \multicolumn{3}{c}{\textbf{What retrieval can do}} \\
\cmidrule(lr){2-4}\cmidrule(lr){5-7}
\textbf{Method} & $N$-ary facts & Typed fact relations & No-shared-entity links & Relation-guided hops & Grounded context & Single-pass query \\
\midrule
MS GraphRAG~\cite{edge2024local}           & \no  & \no  & \no             & \no             & \half$^{a}$ & \yes \\
HippoRAG~2~\cite{gutierrez2025rag}         & \no  & \no  & \no             & \no             & \yes        & \yes \\
MemGraphRAG~\cite{wu2026memgraphrag}       & \no  & \no  & \no             & \no             & \yes        & \yes \\
HyperGraphRAG~\cite{luo2025hypergraphrag}  & \yes & \no  & \no             & \no             & \half$^{a}$ & \yes \\
PropRAG~\cite{wang2025proprag}             & \yes & \no  & \no             & \no             & \yes        & \yes \\
HyGRAG~\cite{zhong2026hygrag}              & \no  & \no  & \half$^{b}$     & \no             & \half$^{a}$ & \yes \\
EventRAG~\cite{yang2025eventrag}           & \yes & \yes & \half$^{c}$     & \half$^{d}$     & \half$^{e}$ & \no  \\
Disco-RAG~\cite{liu2026discorag}           & \no  & \half$^{f}$ & \half$^{f}$ & \no          & \half$^{a}$ & \no  \\
LogicRAG~\cite{chen2026logicrag}           & \na  & \na  & \na             & \na             & \no         & \no  \\
\midrule
\textbf{\sys{} (ours)}                     & \yes & \yes & \yes            & \yes            & \yes        & \yes \\
\bottomrule
\end{tabular}
\par\smallskip
\begin{minipage}{0.97\textwidth}
\footnotesize
$^{a}$Corpus text mixed with LLM-written entity descriptions, community reports, summaries, or plans. $^{b}$Only inside generated cluster summaries. $^{c}$Relations are extracted within one chunk; cross-document links come from shared participants or from the LLM's own knowledge. $^{d}$An LLM agent decides each step over several calls. $^{e}$The graph may contain knowledge the LLM adds from its own parameters. $^{f}$Typed relations between passages, built per question over the passages already retrieved.
\end{minipage}
\end{table*}
\abid{Check the MemGraphRAG and PropRAG rows against their retrieval code (number of query-time LLM calls), and decide whether to add CausalRAG and LightRAG rows.}


% =============================================================================
\section{Methodology}
\label{sec:method}
% =============================================================================

Given a corpus $\mathcal{D} = \{d_1, \ldots, d_n\}$ of documents and a question, \sys{} returns a small set of evidence chains extracted from $\mathcal{D}$, from which a reader LLM answers the question. The work is split into two phases. An offline phase builds the index $\mathcal{G}$ once per corpus (Section~\ref{sec:construction}, Figure~\ref{fig:construction}). An online phase retrieves from $\mathcal{G}$ for each question (Section~\ref{sec:retrieval}, Figure~\ref{fig:retrieval}). The text follows the policy example of Figure~\ref{fig:motivation} throughout. Figures~\ref{fig:construction} and~\ref{fig:retrieval} illustrate the same steps on a small open-domain example. Below, $\simcos(a, b)$ is the cosine similarity between the embeddings of two texts $a$ and $b$; a set of texts is embedded as one concatenated string.

% ─────────────────────────────────────────────────────────────────────────────
\subsection{Index Structure}
\label{sec:index}
% ─────────────────────────────────────────────────────────────────────────────

\paragraph{\textbf{Graph structure.}}
A \sys{} index is a graph
\[
  \mathcal{G} = (C, F, V, \Ecf, \Efv, \Eff)
\]
with three kinds of nodes and three kinds of edges (Figure~\ref{fig:motivation}e).

\textit{Chunks} $C$ are the source passages of the corpus. Each carries its text and a one-sentence summary.

\textit{Facts} $F$ are atomic, self-contained propositions stated in a chunk. Each fact is a tuple
\[
  f = \bigl(\,\underbrace{s}_{\text{sentence}},\;
             \underbrace{T}_{\text{topics}},\;
             \underbrace{A}_{\text{answers}},\;
             \underbrace{U}_{\text{follow-ups}}\,\bigr),
\]
where $s$ is the fact's sentence with every reference resolved, $T$ is a set of topics, $A$ is the set of questions the fact answers, and $U$ is the set of follow-up questions a reader of the fact would ask next. For the grace-period clause in Figure~\ref{fig:motivation}, $f_3$ has $s =$ ``The policy remains in effect during a 30-day grace period after a premium is due and unpaid'', $T = \{$premium payment, grace period$\}$, $A = \{$What happens when a premium is not paid by its due date?$\}$, and $U = \{$What happens if the premium is still unpaid when the grace period ends?$\}$.
\abid{Claude: Bad example the follow and answered question are the exact same. We need a running example or story that is perfectly engineered for UNH and then decide what is the perfect outputs from UNH to make this example work.}

\textit{Entities} $V$ each carry a canonical name $n$ and a \emph{micro-role} $\rho$, a short phrase stating what the entity is or does in its passage. For example, \emph{insured person} has the role ``person whose life is insured under the policy''.

The three edge sets record provenance, participation, and dependence. The provenance relation $\Ecf \subseteq C \times F$ links each fact to the chunks that state it. The participation relation $\Efv \subseteq F \times V$ links each fact to the entities that take part in it. Its participant set $V(f)$ makes each fact a hyperedge of any arity over $V$. The relation set $\Eff$ holds directed, labelled edges $(f_i, f_j, \ell, r)$ between facts. The label $\ell$ is a short type naming how the two propositions relate (\rel{causes}, \rel{precedes}, \rel{exception\_to}, \rel{contradicts}), and $r \in [0, 1]$ is the strength the LLM assigns to the relation. Edges come in pairs that name both directions, such as $(f_2, f_1, \rel{exception\_to})$ and $(f_1, f_2, \rel{has\_exception})$.

\paragraph{\textbf{The dual role of facts.}}
Facts are the only nodes that are connected both to entities and to each other. Through $\Efv$, a fact is a hyperedge: it keeps together all participants of one event, which a triple graph would split across several edges. Through $\Eff$, a fact is a node of a relation graph: it is linked to the facts it causes, precedes, qualifies, or contradicts, whether or not they share a participant. The two edge families support complementary moves. Participation leads from an entity to every fact it takes part in, and back. Relations lead from a fact to the facts it depends on or that depend on it, and say how. The other components support these moves. Facts are the only evidence shown to the reader. Entities route retrieval. Chunks record provenance and are never traversed. The question fields $A$ and $U$ are used only to decide which facts to compare during construction.
\abid{Claude: Why are saying facts are the only evidence shown to the reader? Don't say this ever again. Entities are shown to the reader and so are edge labels.}

\paragraph{\textbf{Edge-restricted views.}}
Removing one edge family gives a familiar index. Without $\Eff$, what remains is the hypergraph $\mathcal{G}_H = (V, \{V(f) : f \in F\})$, an $n$-ary fact index of the kind used by hypergraph RAG~\cite{luo2025hypergraphrag}. Without $V$ and $\Efv$, what remains is the relation graph $\mathcal{G}_R = (F, \Eff)$, a labelled graph over facts. The full index connects more than either view.

\begin{proposition}[Edge-restricted views]
\label{prop:views}
There exist a \sys{} index $\mathcal{G}$ and facts $f, f'$ such that $f'$ is reachable from $f$ in $\mathcal{G}$ but in neither $\mathcal{G}_H$ nor $\mathcal{G}_R$.
\end{proposition}

\begin{proof}
Let $F = \{f_1, f_2, f_3\}$ and $V = \{v_1, v_2\}$, with $V(f_1) = V(f_2) = \{v_1\}$, $V(f_3) = \{v_2\}$, and one relation edge from $f_2$ to $f_3$. In $\mathcal{G}$, the path $f_1 \to v_1 \to f_2 \to f_3$ reaches $f_3$. In $\mathcal{G}_H$, $f_3$ shares no entity with $f_1$ or $f_2$; in $\mathcal{G}_R$, $f_1$ has no edge. Hence $f_3$ is unreachable from $f_1$ in both views.
\end{proof}

The running example combines both kinds of step on one path. Retrieval can enter at the entity \emph{waiver of premiums}, reach the waiver $f_1$ through participation, and follow relations from $f_1$ to the exception $f_2$ and on to the death benefit $f_5$. In $\mathcal{G}_H$, the only route from $f_2$ to $f_3$ runs through \emph{premium}, where it is indistinguishable from fifty-nine others. In $\mathcal{G}_R$, the path cannot start at an entity. Proposition~\ref{prop:views} also underlies our ablation (Section~\ref{sec:ablation}). Restricting traversal to $\Efv$ or to $\Eff$ gives hypergraph-only and relation-only variants over the \emph{same} facts and entities, so their differences from the full system come from the edges available at traversal time, not from extraction.

% ─────────────────────────────────────────────────────────────────────────────
\subsection{Index Construction}
\label{sec:construction}
% ─────────────────────────────────────────────────────────────────────────────

\begin{figure*}[t]
  \centering
  \figorbox{figures/construction.pdf}{\textwidth}{7cm}
  \caption{Index construction, illustrated on a small open-domain example: heavy rain, a flood, and residents leaving, reported by different sources. Steps 1--4 run in parallel for every chunk and produce a chunk subgraph of facts, entities, and relations. Steps 5--8 merge the subgraphs, resolve duplicates, and add relations between facts of different chunks. Here, step~8 relates the rain and the flood, which share no entity.}
  \label{fig:construction}
\end{figure*}

Construction turns $\mathcal{D}$ into $\mathcal{G}$ in eight steps (Figure~\ref{fig:construction}). Steps 1--4 run independently for every chunk. Steps 5--8 work on the merged graph. Every LLM call fills a typed output record~\cite{khattab2024dspy}, without chain-of-thought. Of the LLM's outputs, only fact sentences and relation labels ever reach the reader. Short passages (HotpotQA, MuSiQue) are one chunk each, prefixed with their title; long documents (QuALITY) are split into overlapping token windows.

\subsubsection*{Step 1: Chunk summary and entities.}
One LLM call per chunk returns a one-sentence summary of the chunk and its entity inventory: canonical names, with pronouns and aliases resolved, each with a micro-role. The summary is later shown to the LLM as context when facts from different chunks are compared, and is never shown to the reader. The inventory fixes the chunk's entities before any fact is extracted, so all facts of a chunk refer to the same entity objects. For the section on unpaid premiums, the inventory is \{\emph{premium}: ``monthly amount due for the coverage''; \emph{grace period}: ``30 days after a missed due date during which the policy stays in effect''; \emph{policy}: ``the term-life contract''\}.

\subsubsection*{Step 2: Atomic facts.}
A second call extracts the chunk's facts as sentences that are \emph{atomic} (one claim, event, rule, or state), \emph{self-contained} (every reference resolved), and \emph{grounded} (stated in the chunk). Self-containment is what lets the reader see facts instead of passages~\cite{chen2024dense,min2023factscore}: ``if not paid by then, it will lapse'' becomes ``If outstanding premiums are not paid by the end of the 30-day grace period, the policy lapses.''

\subsubsection*{Step 3: Fact annotation.}
Facts are annotated in batches. Given the chunk, a batch of facts, and the inventory names, one call returns for each fact its topics, the questions it answers, the follow-up questions it raises, and the inventory entities that take part in it. Each returned name is matched to the inventory, exactly or else to the most similar name under character $n$-gram TF--IDF, and becomes a participation edge. Choosing participants from the inventory, rather than extracting names afresh, keeps a chunk's facts attached to shared entity objects.

The answered and follow-up questions are what make relations across sections and documents findable. They describe what a fact is \emph{for} rather than what it \emph{says}. The drunk-driving exception $f_2$ raises ``What happens if premiums are not waived and go unpaid?''. The grace-period clause $f_3$, pages away, answers ``What happens when a premium is not paid by its due date?''. The two sentences share only the hub entity \emph{premium}, but these questions match closely.

\subsubsection*{Step 4: Relations within a chunk.}
For each chunk with at least three facts, one call receives the chunk summary and the numbered facts. It returns every related pair with a forward label, a backward label, and a strength in $[0, 1]$. Labels are constrained in two ways. First, a label must name a relation \emph{between} two propositions (\rel{causes}, \rel{precedes}, \rel{exception\_to}, \rel{contradicts}), not an attribute of one of them (\rel{has\_amount}), which would duplicate participation. Second, the backward label must name the inverse (\rel{causes}/\rel{caused\_by}). Pairs with zero strength are discarded. Triple-based systems store a free-text description on each relation~\cite{edge2024local,guo2024lightrag}. A \sys{} edge stores only its label and strength: the facts already carry the content, and the label is enough to steer traversal and to tell the reader how two facts connect.

\subsubsection*{Step 5: Entity resolution and fact consolidation.}
The chunk subgraphs are merged, and duplicate entities are resolved in passes of increasing cost. The first pass is exact match after normalization (Unicode NFKC, lower-casing, singularization). The second is semantic hashing at high similarity.\footnote{SemHash, \url{https://github.com/MinishLab/semhash}.} The third handles the remaining entities with a nearest-neighbour graph that weighs name similarity above role similarity. Leiden~\cite{traag2019leiden} partitions this graph, and small clusters are packed together. One LLM call per cluster then decides which members denote the same entity and assigns a canonical name and role; edges of duplicates are redirected to the canonical node. The micro-role is what keeps namesakes apart. ``Jordan'' the river and ``Jordan'' the basketball player have identical names but dissimilar roles, while ``NYC'', ``New York'', and ``New York City'' merge. Duplicate facts stated by several chunks are then merged the same way, over a graph of fact similarity. The result is one fact whose provenance is the union of its sources, following the conflict taxonomy of MemGraphRAG~\cite{wu2026memgraphrag}.
\abid{Fact consolidation (\texttt{fact\_dedup}) defaults to on in \texttt{evaluation\_scripts/unh/main.py}, but NOTES.md says it was switched off for the runs. Confirm which setting the reported runs used, and delete the last two sentences if it was off.}

\subsubsection*{Step 6: Multi-view fact affinity.}
Relations between facts of different chunks supply the connectivity that shared entities cannot. Asking an LLM about all pairs, however, is quadratic. Our design choice is to score candidate pairs with several cheap signals, each capturing a different reason two facts may be related, and to send only promising groups to the LLM. The affinity of $f_j$ as a partner of $f_i$ is the sum of four similarities:
\[
  \mathrm{aff}(f_i, f_j) =
  \underbrace{\simcos(s_i, s_j)}_{\text{sentence}} +
  \underbrace{\simcos(T_i, T_j)}_{\text{topics}} +
  \underbrace{\simcos(R_i, R_j)}_{\text{roles}} +
  \underbrace{\simcos(U_i, A_j)}_{\text{raises}\,\to\,\text{answers}},
\]
where $R_i$ is the set of micro-roles of $f_i$'s participants. The first three signals find facts that are about the same thing. The fourth is asymmetric and finds facts that complete each other. It compares the follow-up questions $f_i$ raises with the questions $f_j$ answers. It is therefore high when $f_j$ supplies what a reader of $f_i$ would look for next, which is the bridge of a multi-hop question, even when the two sentences share no words or entities beyond a hub, as with the exception and the grace period.

\subsubsection*{Step 7: Fact clustering.}
Each fact keeps as candidates the $k$ facts with the highest affinity to which it is not yet related. The union of these neighbourhoods is a sparse weighted graph with at most $k|F|$ edges. Hierarchical Leiden~\cite{traag2019leiden} partitions it into communities of at most $M$ facts. A community may span chunks and documents. It is a unit of comparison, not of summarization, and no text is generated for it. Because the affinity rewards facts that complete each other, a community can join the drunk-driving exception with the grace period, which a cluster built from shared entities would keep apart.

\subsubsection*{Step 8: Relations within each community.}
For each community, one call receives its facts as a numbered list. Each fact is preceded by the summary of its chunk, for example ``\texttt{[Context:} exceptions to the waiver of premiums\texttt{]} Premiums are not waived if the disability results from an accident while the insured's blood alcohol exceeds the legal limit''. The call returns the related pairs with labels and strengths under the constraints of Step~4. The summary tells apart similar sentences about different subjects and is not stored. Judging communities rather than pairs has three effects. The LLM sees each candidate alongside its other likely partners and can choose among them. It can relate facts that were grouped transitively rather than directly. The number of calls is bounded by the number of communities, which grows linearly with the number of facts. In the running example, this step adds $f_2$ \rel{leads\_to} $f_3$ between two sections that share only the hub \emph{premium}.

This step differs from community summarization~\cite{edge2024local,zhong2026hygrag} in its output, not in its input. A summary writes the relations among a community's facts into one paragraph. Step~8 writes each relation as a separate labelled edge between two facts. Edges from different communities meet at shared facts and entities, so retrieval can chain them across community boundaries. They can be selected one at a time, so a question receives only the relations on its path. And they leave the facts unchanged, so the reader still sees each clause with its conditions.

\paragraph{\textbf{Cost.}}
Steps 1, 2, and 4 issue one call per chunk, step~3 one per batch of facts, step~5 one per cluster of candidate duplicates, and step~8 one per community. Every term is linear in the number of chunks or facts, so construction cost grows linearly with the corpus (Section~\ref{sec:efficiency}).

% ─────────────────────────────────────────────────────────────────────────────
\subsection{Retrieval}
\label{sec:retrieval}
% ─────────────────────────────────────────────────────────────────────────────

\begin{figure*}[t]
  \centering
  \figorbox{figures/retrieval.pdf}{\textwidth}{6cm}
  \caption{Retrieval, illustrated on the open-domain example of Figure~\ref{fig:construction} for ``\emph{Why did residents leave the old town?}''. One LLM call plans the query. Facts and entities are scored in one pool and diverse seeds are selected. A beam search follows the relation labels the plan asks for: from the seed \emph{residents} it reaches $f_3$, then $f_2$ by \rel{preceded\_by}, then $f_1$ by \rel{caused\_by}, although $f_1$ shares no entity with either. Paths are ranked by a cross-encoder, and the reader receives deduplicated fact chains.}
  \label{fig:retrieval}
\end{figure*}

Retrieval turns a question into a small set of fact chains in six steps (Figure~\ref{fig:retrieval}). Only the first calls an LLM; every later step compares embeddings or applies a cross-encoder. We write $\best(P, b) = \max_{a \in P} \simcos(a, b)$ for the best match between a set of plan items $P$ and a text $b$, and $\cover(P, B) = \frac{1}{|P|}\sum_{a \in P} \max_{b \in B} \simcos(a, b)$ for how well a set of texts $B$ covers the plan items.

\subsubsection*{Step 1: Query planning.}
One LLM call turns the question into a plan with five fields. They are the topics $P_T$, the named entities $P_N$, entity roles $P_R$ (e.g., ``person whose life is insured''), the relation labels $P_L$ the answer path should follow, and short phrases $P_S$ describing the connection sought. For the running example, the plan is $P_T = \{$waiver of premiums, lapse, death benefit$\}$, $P_N = \{$waiver of premiums$\}$, $P_R = \{$person whose life is insured$\}$, $P_L = \{\rel{exception\_to}, \rel{leads\_to}, \rel{prevents}\}$, and $P_S = \{$whether coverage stays in force$\}$. Each field holds one to three items and is left empty rather than filled with invented content. The question itself is not rewritten. The plan is the only point at which an LLM interprets the question.

\subsubsection*{Step 2: Joint scoring of facts and entities.}
Facts and entities compete in one candidate pool:
\begin{align*}
  \mathrm{score}(f) &= w_s\, \simcos(\text{question}, s_f) + w_T\, \cover(P_T, T_f), \\
  \mathrm{score}(v) &= w_n\, \best(P_N, n_v) + w_\rho\, \best(P_R, \rho_v).
\end{align*}
When the plan or the fact has no topics, $w_T$ moves to the sentence term. A question that names its subject enters the index through the entity, here \emph{waiver of premiums}. A question that describes a situation without naming its participants enters through the fact, here the exception $f_2$, whose sentence matches ``driving over the blood-alcohol limit''.

\subsubsection*{Step 3: Diverse seeds.}
Seeds are selected from the pool by Maximal Marginal Relevance~\cite{carbonell1998use}, which adds at each step the candidate $u$ that maximizes
$\lambda\, \mathrm{score}(u) - (1 - \lambda) \max_{u' \in \mathcal{S}} \simcos(u, u')$,
where $\mathcal{S}$ is the current seed set. A fact is represented by its sentence and an entity by its role. Without this step, the seeds of a multi-hop question tend to be paraphrases of one fact. In the running example, the waiver's four other exceptions (attempted suicide, war, crime, and drug misuse) closely resemble $f_2$ and would otherwise crowd out the entity \emph{waiver of premiums}.

\subsubsection*{Step 4: Label-guided beam search.}
From each seed, paths expand over participation edges in both directions and over relation edges, never revisiting a node. Each step is scored by signals that depend on its edge type:
\begin{itemize}
  \item \textbf{Relation step} $f_i \to f_j$ with label $\ell$ and strength $r$: $\alpha_L \cover(P_L, \{\ell\}) + \alpha_S \cover(P_S, \{\ell\}) + \alpha_r\, r$. The step is preferred when its label matches a relation or connection in the plan and the LLM judged it strong.
  \item \textbf{Fact-to-entity step} $f \to v$: $\cover(P_R, \{\rho_v\})$. The step is preferred when the entity's role matches a role in the plan.
  \item \textbf{Entity-to-fact step} $v \to f$: participation edges carry no label in the evaluated configuration. The choice among an entity's facts is therefore left to the diversity term and to the cross-encoder.
\end{itemize}
Each path keeps its $B$ best successors under MMR, using label embeddings for diversity, and ends after $D$ facts, after $2D$ steps, or when no unvisited successor remains. In the running example, from the seed $f_2$, the search prefers \rel{leads\_to} to the grace period $f_3$ over the sixty facts reachable through \emph{premium}, because the label matches the plan. It then follows \rel{precedes} to the lapse rule $f_4$ and \rel{prevents} to the death benefit $f_5$, which an entity hop could reach only through the hub \emph{policy}, shared by 88 sentences. A second path from the same seed follows \rel{exception\_to} to the waiver $f_1$.

\subsubsection*{Step 5: Path ranking.}
Each completed path is serialized as alternating nodes and edges, for example
\begin{center}
\small\ttfamily
"No waiver if disability results from drunk driving"\\
-> LEADS\_TO -> "30-day grace period for unpaid premiums"\\
-> PRECEDES -> "Policy lapses if premiums stay unpaid"
\end{center}
and scored against the question by a cross-encoder~\cite{nogueira2019passage}. A final MMR pass over path embeddings keeps $P$ paths. The cross-encoder reads the whole chain at once, so it judges whether the chain, not only its last fact, answers the question.

\subsubsection*{Step 6: Context assembly.}
Paths kept for one question often share prefixes. The best path is emitted in full. Each remaining path that overlaps it contributes only its steps that have not yet been emitted, and a path that overlaps none starts a new group. The context is thus a set of chains of extracted fact sentences joined by relation labels and by entities with their roles. Each fact appears once and can be traced to its chunk. For the running example, the context contains the chain $f_1$ \rel{has\_exception} $f_2$ \rel{leads\_to} $f_3$ \rel{precedes} $f_4$ \rel{prevents} $f_5$, each fact written out as a sentence. From this chain the reader can conclude that the benefit is not paid. It can also see why: the exception removed the waiver, not the coverage.

\paragraph{\textbf{Cost.}}
Retrieval issues one LLM call per question, for the plan. The rest consists of embedding lookups, a beam search over a sparse graph, and one cross-encoder pass over the candidate paths. Hyperparameters are listed in Appendix~\ref{app:hyper}.


% =============================================================================
\section{Experiments}
\label{sec:experiments}
% =============================================================================

We address the following research questions:
\begin{itemize}
    \item[\textbf{RQ1}] Does \sys{} answer more accurately than existing graph-based RAG methods under a controlled protocol?
    \item[\textbf{RQ2}] Do typed relations between facts add to entity participation, are the two edge families complementary, and do edges beat summaries written over the same clusters?
    \item[\textbf{RQ3}] Where do relations help most: in long narratives, where facts are linked by cause and time, or in encyclopedic multi-hop questions, where they are linked by names?
    \item[\textbf{RQ4}] How much context does \sys{} pass to the reader, and at what construction cost, relative to existing systems at their default settings?
\end{itemize}

% ─────────────────────────────────────────────
\subsection{Experimental Setup}
\label{sec:setup}

\paragraph{Datasets.}
For multi-hop reasoning we use the HotpotQA~\cite{yang2018hotpotqa} and MuSiQue~\cite{trivedi2022musique} releases of HippoRAG~2~\cite{gutierrez2025rag}, also used by MemGraphRAG~\cite{wu2026memgraphrag}, LogicRAG~\cite{chen2026logicrag}, and LinearRAG~\cite{zhuang2025linearrag}: 1{,}000 questions each over pooled corpora of 9{,}811 and 11{,}656 Wikipedia passages. All systems are compared on a fixed, seeded 250-question subset of each, stratified by question type (HotpotQA: 203 bridge, 47 comparison) or hop count (MuSiQue: 129 two-, 79 three-, 42 four-hop); full-set baseline results are in Appendix~\ref{app:full}. For reading comprehension we use QuALITY~\cite{pang2022quality}: multiple-choice questions about stories and articles of about 5k tokens, whose HARD subset cannot be answered by skimming. Its answers hinge on causal, temporal, and motivational links between events rather than on distinct names. Following HyGRAG~\cite{zhong2026hygrag}, each system indexes the 115 articles of the official development split (0.66M tokens) in one pooled index and answers all 2{,}086 questions (1{,}065 HARD).
\abid{Optional: add the insurance-contract benchmark (5 issued policies, 445 questions in five categories) as a further column group of Table~\ref{tab:main} once it is run.}

\paragraph{Corpora.}
The baselines index the full corpora. \sys{} currently indexes only the passages of the 250 evaluated questions (2{,}488 HotpotQA and 4{,}025 MuSiQue passages, the standard distractor setting), so its rows are not yet strictly comparable to the baseline rows.
\abid{Before submission, either build UNH over the full corpora (\texttt{UNH\_CORPUS=full}) or rebuild the baselines on the subset corpora, and then delete this paragraph's caveat.}

\paragraph{Baselines.}
We compare with \textbf{HippoRAG~2}~\cite{gutierrez2025rag} (triples and passages, Personalized PageRank), \textbf{LogicRAG}~\cite{chen2026logicrag} (no pre-built graph; query-time sub-question DAG with a rolling summary), \textbf{MemGraphRAG}~\cite{wu2026memgraphrag} (memory-guided multi-agent triple extraction), \textbf{HyperGraphRAG}~\cite{luo2025hypergraphrag} ($n$-ary hyperedges), \textbf{Microsoft GraphRAG}~\cite{edge2024local} (version 3.2.0, local search), and \textbf{HyGRAG}~\cite{zhong2026hygrag} (hybrid chunk--entity hierarchy with fused summaries). Each system runs from its official code with its default settings, ingestion, and retrieval; we change only model endpoints. The reader receives the context each system's own answer step would use: the top-5 passages (HippoRAG~2), the top-5 256-token windows (MemGraphRAG), the final rolling summary after at most three rounds of top-5 retrieval (LogicRAG), the hybrid entity/hyperedge/source context (HyperGraphRAG), the local-search context with a 12k-token budget (GraphRAG), and one community summary with four chunks, five entities, and five relations (HyGRAG). On QuALITY, the two systems without a chunker of their own (HippoRAG~2 and LogicRAG) receive the 1{,}200/100-token chunks produced by HyGRAG's chunker; the others chunk the articles themselves.

\paragraph{Controlled protocol and metrics.}
Every system uses the same LLM for all internal calls (DeepSeek~V4.1 Flash, reasoning mode off, provider-default temperature), the same embedding model (voyage-4-nano), and the same reader: a shared answer step that applies HippoRAG~2's one-shot QA prompt to the system's context (for QuALITY, a prompt that returns the option letter). Only the context differs across systems. We report exact match (EM) and token-level F1 over the gold answer and its aliases, and LLM-judged accuracy (LLM-Acc) with the judging prompt of LinearRAG~\cite{zhuang2025linearrag}, where the judge sees only the predicted and gold answers; for QuALITY we report the accuracy of the chosen option overall and on HARD questions. \sys{} is compared with each baseline on the same questions by a paired bootstrap (10{,}000 resamples, 95\% intervals); at $n = 250$ one question is 0.4 points and a 95\% interval on F1 spans roughly $\pm 4$--$5$ points. \sys{} uses \texttt{ms-marco-MiniLM-L6-v2} as its cross-encoder and the hyperparameters of Table~\ref{tab:hyperparams}; like the baselines, its context is not truncated.
\abid{Confirm UNH chunking for QuALITY (600/50-token windows by default; the baselines received 1200/100-token chunks).}

% ─────────────────────────────────────────────
\subsection{Main Results (RQ1)}
\label{sec:main}

\begin{table*}[t]
\centering
\caption{Overall QA results. All systems share the LLM, embedding model, reader prompt, and judge. HotpotQA and MuSiQue: 250-question subsets; QuALITY: development set (HARD: hard subset). LLM: LLM-judged accuracy. \sys{}-H and \sys{}-R restrict traversal to participation edges and to relation edges, respectively, on the same index as \sys{}. \sys{}-Sum replaces the relation edges of Step~8 with one relation-focused summary per community. Bold: best; underline: second best (over the rows filled so far). QuALITY: LogicRAG and MS GraphRAG still running.}
\label{tab:main}
\small
\setlength{\tabcolsep}{2pt}
\renewcommand{\arraystretch}{1.1}
\begin{tabular}{lC{1.3cm}C{1.3cm}*{6}{C{1.1cm}}}
\toprule
\multirow{3}{*}{\textbf{Method}} & \multicolumn{2}{c}{\textbf{Reading Comp.}} & \multicolumn{6}{c}{\textbf{Multi-Hop Reasoning}} \\
\cmidrule(lr){2-3}\cmidrule(lr){4-9}
 & \multicolumn{2}{c}{QuALITY} & \multicolumn{3}{c}{HotpotQA} & \multicolumn{3}{c}{MuSiQue} \\
\cmidrule(lr){2-3}\cmidrule(lr){4-6}\cmidrule(lr){7-9}
 & Acc & HARD & EM & F1 & LLM & EM & F1 & LLM \\
\midrule
HippoRAG 2      & \textbf{82.7} & \textbf{76.6} & \textbf{59.6} & \underline{71.9} & 80.4 & 36.0 & 47.1 & 52.0 \\
LogicRAG        & \tbd & \tbd & \underline{56.4} & 71.3 & \underline{81.6} & 33.6 & 44.5 & 48.8 \\
MemGraphRAG     & 74.8 & 67.0 & 50.4 & 61.4 & 68.4 & 30.4 & 38.8 & 41.6 \\
HyperGraphRAG   & 79.4 & 71.7 & \textbf{59.6} & \textbf{74.8} & \textbf{85.6} & \textbf{46.4} & \textbf{57.2} & \textbf{62.0} \\
MS GraphRAG     & \tbd & \tbd & 51.6 & 65.4 & 75.6 & \underline{38.8} & \underline{49.7} & \underline{55.2} \\
HyGRAG          & \underline{81.7} & \underline{75.8} & \underline{56.4} & 69.3 & 79.2 & 32.0 & 40.4 & 44.0 \\
\midrule
\sys{}-H (hypergraph) & \tbd & \tbd & \tbd & \tbd & \tbd & \tbd & \tbd & \tbd \\
\sys{}-Sum (summaries) & \tbd & \tbd & \tbd & \tbd & \tbd & \tbd & \tbd & \tbd \\
\sys{}-R (relations) & \tbd & \tbd & \tbd & \tbd & \tbd & \tbd & \tbd & \tbd \\
\midrule
\textbf{\sys{} (ours)} & \tbd & \tbd & \tbd & \tbd & \tbd & \tbd & \tbd & \tbd \\
\bottomrule
\end{tabular}
\end{table*}

Table~\ref{tab:main} reports the results. Under the controlled protocol, HyperGraphRAG is the strongest baseline on both multi-hop datasets (tied with HippoRAG~2 on HotpotQA EM), and its margin is largest on MuSiQue (57.2 F1, 7.5 points above the next system), consistent with the argument that preserving $n$-ary facts helps when evidence must be combined. It also passes by far the longest context of any system to the reader (about 21.7k tokens; Section~\ref{sec:efficiency}). HippoRAG~2 and LogicRAG follow on HotpotQA, and Microsoft GraphRAG on MuSiQue. MemGraphRAG ranks lower than in its own paper, which used a different LLM, embedding model, and metric, illustrating why a shared protocol is needed before published numbers are compared. Baseline rankings on the subsets match those on all 1{,}000 questions up to one adjacent swap per dataset (Appendix~\ref{app:full}).
\abid{Fill the UNH rows and paired-bootstrap intervals. Discuss: margin over HyperGraphRAG at a fraction of its context; whether the gain grows with hop count on MuSiQue; failures where seeds were correct but traversal was not.}

% ─────────────────────────────────────────────
\subsection{Relations versus Entities and Summaries (RQ2)}
\label{sec:ablation}

To isolate the contribution of the relation graph, we evaluate two views of the \emph{same} index (Proposition~\ref{prop:views}) with the same seeds, ranking, and reader. \textbf{\sys{}-H} traverses only participation edges $\Efv$, so facts reach each other only through shared entities, as in hypergraph RAG. \textbf{\sys{}-R} traverses only relation edges $\Eff$; entities can still serve as seeds but cannot be traversed through. Because extraction and entity resolution are shared, the differences among \sys{}-H, \sys{}-R, and \sys{} (Table~\ref{tab:main}) come from the edges available at traversal time. \sys{}-R versus \sys{}-H measures whether typed relations between facts retrieve better evidence than entity co-membership. \sys{} versus both measures whether the two edge families are complementary.
\abid{Fill and discuss; report the share of retrieved paths that mix both edge families, and the share of relation steps between facts that share no entity.}

\paragraph{Edges versus summaries.}
A summary of a community can, in principle, state the same relations that Step~8 stores as edges, and a prompt can ask it to. \textbf{\sys{}-Sum} tests this alternative in its strongest form. It keeps \sys{}'s facts, entities, and communities, so the grouping is the same, including the question-based affinity that already joins facts which an entity-based clustering would separate. In place of Step~8 it asks the LLM, for each community, for one summary that states every causal, temporal, conditional, and exception relation among the community's facts. At query time, the seeds are scored as in \sys{}, and the reader receives the summaries of the seeds' communities together with the seed facts. \sys{} versus \sys{}-Sum therefore isolates the form in which relations are stored: edges that retrieval selects and chains per question, or text written once per community. Microsoft GraphRAG and HyGRAG in Table~\ref{tab:main} add the comparison with summaries over entity-based clusters.
\abid{Implement \sys{}-Sum (same index; community summary prompt; summaries + seed facts as context). Report accuracy, reader-context length, and construction tokens against \sys{}. Also count, on a sample of questions, how often the gold chain crosses two or more communities, since a single summary cannot cover such chains.}

\paragraph{Case study.}
\abid{Trace one QuALITY HARD question end to end, in the style of Figure~\ref{fig:motivation}: the plan, the seeds, the path \sys{} follows (with its relation labels and source chunks), and the context of the strongest baseline for the same question. Pick a question whose path crosses a relation between facts that share no entity.}

% ─────────────────────────────────────────────
\subsection{Effectiveness across Domains (RQ3)}
\label{sec:domains}

Wikipedia multi-hop questions reward connecting short encyclopedic passages through named entities, which is the setting entity-centric indices are designed for. QuALITY tests the opposite regime: long narratives in which relevant facts are linked by causation, temporal order, and motivation more often than by distinct names. We expect the relation graph to matter most on QuALITY's HARD questions, which require integrating evidence from several parts of a story. Among the baselines, HippoRAG~2 (82.7) and HyGRAG (81.7) lead, within about one standard error of each other ($\approx$0.9 points at $n = 2{,}086$), followed by HyperGraphRAG (79.4) despite a context four times longer, and MemGraphRAG (74.8), whose 256-token windows give the shortest context. HARD questions keep this order and cost every system 6--8 points.
\abid{Fill LogicRAG, MS GraphRAG, and \sys{} on QuALITY and discuss; add further domains as column groups.}

% ─────────────────────────────────────────────
\subsection{Reader Context and Construction Cost (RQ4)}
\label{sec:efficiency}

Table~\ref{tab:cost} reports the LLM tokens consumed by index construction and the mean length of the context passed to the reader on the evaluated questions. Construction is paid once per corpus; context is paid on every query.

\begin{table*}[t]
\centering
\caption{Construction cost and reader context. In/Out: construction LLM input and output tokens (millions); Ctx: mean reader-context tokens per question (\texttt{cl100k\_base}; context only). Baselines index the full HotpotQA and MuSiQue corpora and the 115 QuALITY development articles; \sys{} currently indexes the subset corpora (Section~\ref{sec:setup}). $^{\dagger}$LogicRAG builds no index; its context is a summary written by the LLM at query time.}
\label{tab:cost}
\footnotesize
\setlength{\tabcolsep}{2pt}
\renewcommand{\arraystretch}{1.15}
\begin{tabular}{l*{9}{C{1.1cm}}}
\toprule
\multirow{2}{*}{\textbf{Method}} & \multicolumn{3}{c}{\textbf{HotpotQA}} & \multicolumn{3}{c}{\textbf{MuSiQue}} & \multicolumn{3}{c}{\textbf{QuALITY}} \\
\cmidrule(lr){2-4}\cmidrule(lr){5-7}\cmidrule(lr){8-10}
 & In & Out & Ctx & In & Out & Ctx & In & Out & Ctx \\
\midrule
HippoRAG 2           & 9.25   & 3.59  & 634      & 10.50  & 3.75  & 664      & 1.86 & 0.39 & 5{,}328 \\
LogicRAG$^{\dagger}$ & 0      & 0     & 66       & 0      & 0     & 115      & 0    & 0    & \tbd \\
MemGraphRAG          & 158.70 & 16.51 & 1{,}241  & 153.45 & 16.05 & 1{,}240  & 45.50 & 4.22 & 1{,}292 \\
HyperGraphRAG        & 161.93 & 17.57 & 21{,}783 & 189.95 & 19.59 & 21{,}540 & 19.72 & 5.75 & 20{,}492 \\
MS GraphRAG          & 47.37  & 21.40 & 9{,}087  & 49.33  & 22.43 & 8{,}898  & \tbd & \tbd & \tbd \\
HyGRAG               & 13.77  & 4.33  & 1{,}000  & 15.26  & 4.52  & 1{,}018  & 2.93 & 0.48 & 5{,}014 \\
\midrule
\textbf{\sys{} (ours)} & \tbd & \tbd & \tbd & \tbd & \tbd & \tbd & \tbd & \tbd & \tbd \\
\bottomrule
\end{tabular}
\end{table*}

\paragraph{Construction.}
Construction cost varies by more than an order of magnitude. HippoRAG~2 and HyGRAG, which issue one entity and one triple-extraction call per passage (plus community summaries for HyGRAG), use 12.8--19.8M tokens. Microsoft GraphRAG spends about two thirds of its 69--72M tokens on community reports. MemGraphRAG and HyperGraphRAG use 170--210M tokens; for MemGraphRAG, two thirds of the input goes to schema induction and most of the rest to conflict detection. On QuALITY the ordering is the same, but long articles make construction cheaper per corpus token for systems that call the LLM once per chunk: HyperGraphRAG spends 30 input tokens per corpus token (125 on HotpotQA) and HippoRAG~2 2.8 (7.1), while MemGraphRAG, whose 256-token windows do not grow, spends 69 (122).
\abid{Fill UNH from \texttt{construction\_results.json} (per-stage split). The NOTES smoke test measured about 7.9k input / 4.8k output tokens per HotpotQA passage; since UNH indexes fewer passages, also report tokens per passage for all systems.}

\paragraph{Reader context.}
Under their defaults, context lengths range from 634 tokens (HippoRAG~2's five passages) through about 1{,}000--1{,}250 (HyGRAG, MemGraphRAG) to 9k (GraphRAG) and 21.7k (HyperGraphRAG). The strongest baseline in Table~\ref{tab:main} is also the one with the longest context, so accuracy and context length are hard to separate across systems. QuALITY preserves this ordering (MemGraphRAG 1.3k, HyGRAG 5.0k, HippoRAG~2 5.3k with its longer chunks, HyperGraphRAG 20.5k tokens). \sys{} returns at most $P = 20$ deduplicated chains of at most $D$ facts each.
\abid{Fill the UNH context length and state the comparison with HyperGraphRAG.}

\paragraph{The case of LogicRAG.}
LogicRAG's final context is the shortest (66 and 115 tokens), but it measures something different. LogicRAG builds no index and does at query time the work other systems do offline: it retrieves five passages, decomposes the question when they do not suffice, retrieves five more passages per sub-problem, and after each round has the LLM rewrite a rolling summary; only the last summary reaches the reader. On the MuSiQue subset it performs 1.3 additional rounds on average and reads about 11.5 passages (about 1{,}270 tokens of raw evidence, counting repeats) to produce its 115-token summary (HotpotQA: 0.4 rounds, 6.9 passages, about 810 tokens), and its query-time pipeline consumes about 2.9k LLM input and 0.6k output tokens per question (HotpotQA: 1.8k and 0.3k, averaged over all 1{,}000 questions, including its own answer and judge calls). Its short context is thus the output of query-time compression rather than of more selective retrieval, and its cost recurs on every question; it should be compared on the evidence it inspects, not on the length of its summary.

\paragraph{Caveat.}
Every system, including \sys{}, runs at its default context size; we did not search for the budget at which each performs best. A system with a long default context might keep most of its accuracy at a smaller budget, so our comparison characterizes the systems as released rather than their best accuracy--context trade-off.

% =============================================================================
\section{Discussion and Limitations}
\label{sec:limitations}
% =============================================================================

\paragraph{Construction cost.}
\sys{} trades construction cost for structure: every relation is adjudicated by an LLM and several calls are issued per chunk. The number of calls is linear in corpus size, but the constant is larger than for extraction-only systems; label-only relation prompts and smaller annotation schemas reduce it.

\paragraph{Inferred relations.}
Facts are extracted from the corpus, but relations between facts of different documents are the LLM's judgment: the policy does not state that a denied waiver leads to a lapse. \sys{} keeps the two apart. A relation stores only a label and a strength, both of its facts remain traceable to their passages, and the reader sees the label as a connector between quoted facts rather than as a claim of the corpus. Wrong relations can still mislead traversal; recording which community call produced each relation would make them auditable.
\abid{Optional: measure relation precision on a sample of 100 cross-chunk edges judged by hand.}

\paragraph{Fact-level consistency.}
Duplicate facts are merged (Step~5), but contradictory facts are linked by a \rel{contradicts} edge rather than reconciled, and the reader must weigh them.
\abid{Reconcile with the fact-consolidation setting the runs used (Step~5 note).}

\paragraph{Open-vocabulary relations.}
Relation labels are produced by the LLM under constraints rather than drawn from a fixed inventory such as the Penn Discourse Treebank's~\cite{prasad2017penn}. Synonymous labels (\rel{causes}, \rel{leads\_to}) are matched only through their embeddings, and participation edges are unlabelled, so expansion out of an entity is not query-guided. Normalizing labels into a shared vocabulary, as ontology induction over \sys{} graphs does~\cite{mridul2026groil}, would also make the index usable as a Semantic Web resource.

\paragraph{Scope.}
\sys{} targets questions whose answer follows a chain of specific statements. It does not target questions about a corpus as a whole (``What are the main themes of these policies?''), for which community summaries are designed~\cite{edge2024local}. The two are complementary: summaries could be written over \sys{} communities and kept alongside the edges. All systems share one LLM and one embedding model, and the multi-hop comparison uses 250 questions per dataset; small differences should be read against the bootstrap intervals.

% =============================================================================
\section{Conclusion}
% =============================================================================

Graph-based RAG indices connect knowledge only where facts mention the same entity, while many questions follow relations between statements: a cause and its effect, a rule and its exception, a claim and its rebuttal. We presented the \sysfull{}, an index in which every atomic fact is both a hyperedge over its participants and a node in a graph of typed relations between facts, each fact traceable to its passage. Making the relation between facts an element of the index lets retrieval move between dependent facts that share no entity, choose its next step by the kind of relation a question asks for, and give the reader compact chains of grounded facts instead of passages or generated summaries. Construction stays linear in corpus size by judging relations within communities of a question-based fact-affinity graph rather than between all pairs. We compared \sys{} with six graph-based RAG systems under a protocol that fixes the language model, embedding model, reader, and judge, isolated the relation graph by comparing views of the same index, and reported every system's construction cost and reader context.
\abid{Add one sentence with the main empirical findings (RQ1--RQ4).}

% ---- Bibliography ----
\bibliographystyle{ACM-Reference-Format}
\bibliography{references}


\appendix
\section{Full-Set Baseline Results}
\label{app:full}

Table~\ref{tab:full} reports the baselines on all 1{,}000 questions of each dataset, with the same reader and judge as Table~\ref{tab:main}. Contain-Acc is the fraction of answers that contain the normalized gold answer (the Str-Acc of MemGraphRAG~\cite{wu2026memgraphrag} and LinearRAG~\cite{zhuang2025linearrag}). Subset F1 lies within 4.2 points of full-set F1 for every system, inside the subset's sampling uncertainty, and the rankings agree except for one adjacent swap per dataset (GraphRAG and HyGRAG on HotpotQA; HippoRAG~2 and GraphRAG on MuSiQue).

\begin{table*}[t]
\centering
\caption{Baselines on the full 1{,}000-question sets (shared reader and judge).}
\label{tab:full}
\small
{%
\begin{tabular}{lcccccccc}
\toprule
\multirow{2}{*}{\textbf{Method}} & \multicolumn{4}{c}{\textbf{HotpotQA}} & \multicolumn{4}{c}{\textbf{MuSiQue}} \\
\cmidrule(lr){2-5}\cmidrule(lr){6-9}
 & EM & F1 & Contain-Acc & LLM-Acc & EM & F1 & Contain-Acc & LLM-Acc \\
\midrule
HippoRAG 2      & 61.3 & 73.5 & 70.6 & 82.8 & 39.8 & 50.6 & 43.9 & 53.9 \\
LogicRAG        & 58.2 & 72.8 & 70.2 & 82.9 & 35.8 & 48.5 & 42.3 & 52.8 \\
MemGraphRAG     & 51.2 & 62.6 & 60.9 & 70.6 & 30.9 & 40.5 & 34.4 & 42.9 \\
HyperGraphRAG   & 61.4 & 74.9 & 74.5 & 86.4 & 44.5 & 56.4 & 51.7 & 62.1 \\
MS GraphRAG     & 56.4 & 69.2 & 67.1 & 79.6 & 38.2 & 49.1 & 42.9 & 54.2 \\
HyGRAG          & 54.1 & 67.5 & 64.8 & 78.6 & 34.1 & 44.5 & 37.1 & 47.7 \\
\bottomrule
\end{tabular}}
\end{table*}


\section{Hyperparameters}
\label{app:hyper}

Table~\ref{tab:hyperparams} lists the values used in all experiments. Symbols follow Section~\ref{sec:method}.

\begin{table}[h]
\centering
\caption{\sys{} hyperparameters used in all experiments.}
\label{tab:hyperparams}
\small
\begin{tabular}{p{1.5cm}p{4.3cm}c}
\toprule
\textbf{Stage} & \textbf{Parameter} & \textbf{Value} \\
\midrule
Construction & facts annotated per call & 15 \\
             & minimum facts for an intra-chunk relation call & 3 \\
             & entity resolution: neighbours / threshold / max.\ cluster & 15 / 0.5 / 72 \\
             & affinity: four views, equal weights & 1, 1, 1, 1 \\
             & candidate neighbours per fact $k$ & 12 \\
             & maximum community size $M$ & 50 \\
\midrule
Seeding      & fact weights $(w_s, w_T)$ & 0.8, 0.2 \\
             & entity weights $(w_n, w_\rho)$ & 0.6, 0.4 \\
             & seeds $K$, MMR $\lambda$ & 12, 0.6 \\
\midrule
Traversal    & edge weights $(\alpha_L, \alpha_S, \alpha_r)$ & 0.3, 0.5, 0.2 \\
             & beam width $B$, MMR $\lambda$ & 6, 0.75 \\
             & maximum facts per path $D$ (steps $2D$) & 4 (8) \\
\midrule
Ranking      & paths kept $P$, MMR $\lambda$ & 20, 0.6 \\
\bottomrule
\end{tabular}
\end{table}
\abid{Values updated to the current run defaults (\texttt{evaluation\_scripts/unh/main.py}): $M$ 40$\to$50, $K$ 9$\to$12, $D$ 3$\to$4, beam MMR 0.6$\to$0.75. Confirm against the final runs, including the seed MMR $\lambda$ and the 15 entity-resolution neighbours, which were not re-checked.}


\end{document}
